%%=================================================================
%% UFAC-modelo-TCC.tex (UTF-8 format), v1.0
%%
%%
%% Adaptação do Modelo de Trabalho de Conclusão de Curso (TCC) do
%% curso de Bacharelado em Sistemas de Informação da Universidade
%% Federal do Acre em LaTeX utilizando o pacote abnTeX2
%% (http://abnTeX2.googlecode.com).
%%
%% Precisa do arquivo abntex2-UFAC.sty e abntex2-alf.bst e
%% utilizando o compilador pdfLaTeX.
%%
%%
%% This work may be distributed and/or modified under the
%% conditions of the LaTeX Project Public License (LPPL), either
%% version 1.3 of this license or (at your option) any later
%% version.
%% The latest version of this license is in
%% http://www.latex-project.org/lppl.txt and version 1.3 or later
%% is part of all distributions of LaTeX version 2005/12/01 or
%% later.
%%
%% Further information about abnTeX2 are available on
%% https://www.abntex.net.br/
%%
%%
%% This work has the LPPL maintenance status `maintained'.
%%
%% The Creator and Current Maintainer of this work:
%% Manoel Limeira <juniorlimeiras@gmail.com>.
%%
%% Others contributors of this work:
%% Christopher Saar <cristophersaar@gmail.com>, Victor Alexandre<victoralexandre922@gmail.com>, Clécio Elias<https://github.com/Cleps>.
%%
%%
%% Última versão janeiro/2024
%%=================================================================

% MODIFIED: Adicionada opção chapter=TITLE para que os títulos dos capítulos, elementos pré-textuais e pós-textuais sejam escritos em maiúsculas
\documentclass[
    % -- opções da classe memoir --
    12pt,				       % Tamanho da fonte
    openright,			       % Capítulos começam em pág. ímpar (insere página vazia caso preciso)
    oneside,			       % Impressão apenas no anverso. Oposto a opção "twoside"
    a4paper,			       % Tamanho do papel
    % -- opções do pacote abntex2 --
    chapter=TITLE,             % Títulos de capítulos, de elementos pré-textuais e pós-textuais em maiúsculas
    sumario=tradicional,       % Sumário padrão memoir. Oposto a opção "abnt-6027-2012"
    % -- opções do pacote babel --
    english,			        % Idioma adicional para hifenização
    brazil, 				    % Idioma principal do documento
 ]{abntex2}


% ADDED: Importação do arquivo abntex2-alf.bst para que o sistema de citação seja (Autor, data) e dos pacotes datetime, xcolor, pgfplots, pgfplotstable e microtype para adicionar funcionalidades de arranjo visual do texto, plotagem de dados, estilização do texto, acesso de datas e horários
%% CONFIGURAÇÕES DE PACOTES

% Pacotes básicos
\usepackage{lmodern}			% Usa a fonte Latin Modern
\usepackage[T1]{fontenc}		% Seleção de códigos de fonte de saída
\usepackage[utf8]{inputenc}		% Codificação do documento (conversão automática dos acentos)
\usepackage{indentfirst}		% Indenta o primeiro parágrafo de cada seção
\usepackage{graphicx}			% Inclusão de gráficos e figuras
\usepackage{booktabs}           % \toprule, \midrule e \bottomrule para tabelas
\usepackage{datetime}           % Acesso de datas e horários na compilação
\usepackage{pgfgantt} % para usar o gráfico de gant

% MODIFIED: Adicionada opção abnt-etal-text=it para que "et al." seja escrito em itálico
% Sistema (Autor, data) com títulos das referências em negrito e termo latino "et al." escrito em itálico
\usepackage[alf,abnt-emphasize=bf, abnt-etal-text=it]{abntex2cite}
\bibliographystyle{abntex2-alf.bst}
% FUTURE-TODO: Modificar comando \cite para usar sistema (Autor, data) e não (AUTOR, data) sem precisar do arquivo .bst - apagar arquivo .bst e comando acima quando finalizado

% Títulos em fonte Times
\usepackage{mathptmx}
\renewcommand{\ABNTEXchapterfont}{\rmfamily\bfseries}
\usepackage{lmodern}
\usepackage{tikz}
\usetikzlibrary{arrows.meta, positioning, shapes}
% Leitura de figuras .eps, compilação para PDF
\usepackage[outdir=./]{epstopdf}
\usepackage{epsfig}

% \usepackage{svg}                % Leitura de figuras .svg na ferramenta OVERLEAF

\usepackage{abntex2-UFAC}       % Personalização de ESTILO para a Universidade Federal do Acre

% Pacotes adicionais

% Plotagem de dados
\usepackage{pgfplots, pgfplotstable} % mais detalhes veja https://pt.overleaf.com/learn/latex/Pgfplots_package
\pgfplotsset{compat=1.15} % resolve o problema do pacote pgfplots em múltiplas ferramentas de compilação

\usepackage{xcolor}             % Controle das cores

\usepackage{microtype} 			% Para aprimoramento do arranjo visual do texto na ferramenta OVERLEAF

\usepackage{amsmath}            % Pacote para controle de equações matemáticas
\usepackage{float}
%\usepackage{fullpage}
\usepackage{enumitem}
\usepackage{pgf-pie}
% FUTURE-TODO: Modificar comandos relacionados as expressões latinas para serem escritas em itálico

% MODIFIED: Termos latinos para serem escritos em itálico
\renewcommand{\apudname}{\textit{apud}}
\renewcommand{\passimname}{\textit{passim}}
\renewcommand{\Idemname}{\textit{Id.}}
\renewcommand{\Ibidemname}{\textit{Ibid.}}
\renewcommand{\opcitename}{\textit{op.\ cit.}}
\renewcommand{\loccitname}{\textit{loc.\ cit.}}
\renewcommand{\cfcitename}{\textit{Cf.}}
\renewcommand{\etseqname}{\textit{et seq.}}
% ---

% MODIFIED: Textos para serem mais diversificados e datas para serem escritas de forma automática
%% Informações de dados para CAPA e FOLHA DE ROSTO
\titulo{Construção e Avaliação de um Benchmark Público em Português Brasileiro para Classificação Automática de Issues em Repositórios de Software}

\autor{Marcos Vinícius Moraes Costa}

\local{Rio Branco}

\data{\the\year} % Automatiza a inserção do ano atual, para modificação por um ano específico troque "\the\year" pelo ano desejado

\orientador{Catarina de Souza Costa} % Redefinido no abntex2-UFAC.sty para aceitar Instituição (default = Curso de Bacharelado em Sistemas de Informação - Universidade Federal do Acre)



\instituicao{Universidade Federal do Acre}

% Campus ou Centro (opcional)
\campus{Centro de Ciências Exatas e Tecnológicas} % Pacote abntex2-UFAC

\curso{Curso de Bacharelado em Sistemas de Informação} % Pacote abntex2-UFAC



\databanca{\today} % Automatiza a inserção da data atual, para modificação por uma data específica troque "\today" pela data desejada por extenso

% REMOVED: Comando \preambulo por ter sido inutilizado no abntex2-UFAC.sty
% O preambulo deve conter o tipo do trabalho (tese, dissertação, trabalho de conclusão de curso e outros); objetivo (aprovação em disciplina, grau pretendido e outros); nome da instituição a que é submetido e a área de concentração
% \preambulo{Projeto de monografia apresentado como exigência parcial para obtenção do grau de bacharel em Sistemas de Informação da Universidade Federal do Acre.}
% ---

%% Configurações de aparência do PDF final

% informações para o arquivo pdf de saída
\makeatletter

% MODIFIED: Cores de links do documento
\hypersetup{
    % metadados
    pdftitle={\@title},
    pdfauthor={\@author},
    pdfsubject={\imprimirpreambulo},
    pdfcreator={LaTeX with abnTeX2},
    %%% IMPORTANTE: ALTERAR "colorlinks=false" PARA "colorlinks=true" ANTES DE IMPRIMIR A VERSÃO FINAL DO DOCUMENTO
    colorlinks=false,           % false: links em frame; true: links coloridos
    linkcolor=black,            % cor dos links no documento
    citecolor=black,            % cor dos links para a bibliografia
    filecolor=black,            % cor dos links para arquivos
    urlcolor=black,             % cor dos links para sites
    bookmarksdepth=4,           % profundidade do sumário do PDF
}

\makeatother
% ---

\begin{document}
% Retira espaço extra obsoleto entre as frases
\frenchspacing


% ----------------------------------------------------------
% ELEMENTOS PRÉ-TEXTUAIS
% ----------------------------------------------------------
\pretextual


%% Capa
\imprimircapa
% ---

%% Folha de rosto
\imprimirfolhaderosto
% ---


    
 %%Lista de ilustrações
\pdfbookmark[0]{\listfigurename}{lof}
\listoffigures*
\cleardoublepage


% % ADDED: Comando para imprimir lista de quadros
%Lista de quadros
\pdfbookmark[0]{\listofquadrosname}{loq}
\listofquadros*
\cleardoublepage
% % ---

% %% Lista de tabelas
%\pdfbookmark[0]{\listtablename}{lot}
%\listoftables*
%\cleardoublepage
% % ---

% % MODIFIED: Texto
% %% Lista de siglas e abreviaturas (opcional)
% % sintaxe: \item [sigla] Descrição da sigla
% \begin{siglas}
% 	\item[ABNT] Associação Brasileira de Normas Técnicas
% 	\item[UFAC] Universidade Federal do Acre
% \end{siglas}
% %% ---

% % MODIFIED: Exemplos de símbolos
% %% Lista de símbolos (opcional)
% % sintaxe: \item [simbolo] Descrição do símbolo
% \begin{simbolos}
%     \item[$\infty$] Infinito
%     \item[$\Gamma$] Letra grega Gama
%     \item[$\Lambda$] Lambda
%     \item[$\zeta$] Letra grega minúscula zeta
%     \item[$\in$] Pertence
% \end{simbolos}
% %% ---

%% Sumário
\pdfbookmark[0]{\contentsname}{toc}
\tableofcontents*
\cleardoublepage
% ---

%% Inicializado contadores
\newcounter{equationCont} % Contador de equações
% ---

% ----------------------------------------------------------
% ELEMENTOS TEXTUAIS
% ----------------------------------------------------------
\textual


% MODIFIED: Adição de quebras de linha para início de parágrafos corretamente, troca de links quebrados, adição ou ajuste de textos de exemplo, posição da legenda e adição de fonte nas figuras e tabelas nos elementos textuais

% ADDED: Comando \pagestyle{simple} para remoção do cabeçalho no modelo Capítulo <N>. <Nome do Capítulo>
\pagestyle{simple}

% Modifique a estrutura dos capítulos e seções de acordo com a necessidade do seu trabalho

%% Introdução
\chapter{APRESENTAÇÃO}\label{sec:introducao}




% problema da pesquisaaaa AA%
\chapter{PROBLEMA DA PESQUISA}\label{sec:problema}

\section{QUESTÕES DE PESQUISA}\label{sec:questoes-pesquisa}

Com base no problema investigado e nos objetivos estabelecidos, este trabalho busca responder às seguintes questões de pesquisa:

\begin{description}
    \item[\textbf{RQ1.}] Em que medida labels heterogêneos de diferentes repositórios públicos podem ser harmonizados em uma taxonomia comum e reproduzível para classificação de issues em português brasileiro?
    \item[\textbf{RQ2.}] Como se compara o desempenho de classificadores baseados em TF-IDF com uma abordagem baseada em representações contextuais do BERTimbau na classificação das classes harmonizadas do benchmark?
    \item[\textbf{RQ3.}] Quais padrões de erro são observados nos classificadores avaliados e em quais classes e características textuais eles se concentram?
    \item[\textbf{RQ4 -- exploratória.}] Como o desempenho dos classificadores varia quando são avaliados em issues provenientes de repositórios não presentes nos dados de treinamento?
\end{description}

A RQ4 será tratada como questão exploratória, condicionada à composição e ao volume do benchmark obtido, de modo a não comprometer a execução das questões principais do estudo.

\chapter{OBJETIVOS DA PESQUISA}\label{sec:objetivos}
Nesta seção estão estabelecidos os objetivos desta pesquisa, tanto de forma geral quanto de forma específica.

\section{OBJETIVO GERAL}\label{sec:ObjGeral}

Construir e avaliar um benchmark público de issues em português brasileiro para classificação automática de demandas em repositórios de software, comparando modelos clássicos de aprendizado de máquina com uma abordagem baseada em BERTimbau.

\section{OBJETIVOS ESPECÍFICOS}\label{sec:ObjEspec}

\begin{enumerate}
    \item Selecionar repositórios públicos com quantidade relevante de issues em português brasileiro e rótulos adequados à classificação.
    \item Coletar e organizar as issues dos repositórios selecionados por meio da GitHub REST API, garantindo a reprodutibilidade do processo.
    \item Definir uma taxonomia unificada de classes a partir da harmonização dos rótulos utilizados nos diferentes repositórios.
    \item Construir e disponibilizar um benchmark público, documentado e reprodutível de issues em português brasileiro, com rótulos originais e classes harmonizadas.
    \item Implementar e avaliar classificadores baseados em TF-IDF combinados com Regressão Logística, SVM e Random Forest.
    \item Comparar os modelos clássicos com uma abordagem baseada em BERTimbau, utilizando métricas de classificação e análise de erros.
\end{enumerate}







%% Capítulo de Resultados
\chapter{JUSTIFICATIVA DA PESQUISA}\label{sec:justificativa}



\chapter{FUNDAMENTAÇÃO TEÓRICA}\label{sec:fundamentacao}

Nesta seção serão abordados alguns conceitos fundamentais para a realização deste trabalho, tais como: 



\section{Assunto 1}





\section{Assunto 2}



\section{Asssunto 3}






\section{Trabalhos Relacionados}\label{sec:TrabRel}



\chapter{PROCEDIMENTOS METODOLÓGICOS}



\chapter{ESBOÇO DOS CAPÍTULOS E SEÇÕES}\label{sec:cronograma}


\hspace{0em} 1. INTRODUÇÃO

\hspace{2em} 1.1. Problema

\hspace{2em} 1.2. Objetivos

\hspace{4em} 1.2.1. Geral

\hspace{4em} 1.2.2. Específicos

\hspace{2em} 1.3. Justificativa

\hspace{2em} 1.4. Metodologia

\hspace{2em} 1.5. Organização do Trabalho

\hspace{0em} 2. FUNDAMENTAÇÃO TEÓRICA

\hspace{2em} 2.1. 

\hspace{2em} 2.2. 

\hspace{2em} 2.3. 

\hspace{2em} 2.4. 

\hspace{0em} 3. 

\hspace{2em} 3.1. 

\hspace{2em} 3.2. 

\hspace{4em} 3.2.1 

\hspace{4em} 3.2.2 

\hspace{4em} 3.2.3 

\hspace{0em} 4. RESULTADOS

\hspace{0em} 5. CONSIDERAÇÕES FINAIS

\hspace{2em} 5.1. Contribuições

\hspace{2em} 5.1 Limitações do Estudo

\hspace{2em} 5.3 Trabalhos Futuros

\hspace{0em} REFERÊNCIAS


\chapter{CRONOGRAMA}\label{sec:cronograma}

O processo de escrita desta monografia seguirá o cronograma explicitado no Quadro 1.

\begin{quadro}[!ht]
    \centering
    \textbf{\caption{Cronograma de execução da pesquisa.}}
    \label{qua: cronograma}
    
    \begin{tabular}{|p{2.2cm}|l|p{1.2cm}|p{7cm}|} \hline
        % Títulos
        \textbf{Mês} & \textbf{Data} & \textbf{Dias} & \textbf{Atividade} \\ \hline
        % Linhas do cronograma
        Novembro & 03-11-2025 até 17-11-2025 & 15 & Escrita do \textbf{Capítulo 1 - Introdução}. \\ \hline
        Novembro & 18-11-2025 até 30-11-2025 & 13 & Início da escrita do \textbf{Capítulo 2 - Fundamentação Teórica}. \\ \hline
        Dezembro & 01-12-2025 até 15-12-2025 & 15 & Conclusão da escrita do \textbf{Capítulo 2 - Fundamentação Teórica}. \\ \hline
        Dezembro & 16-12-2025 até 23-12-2025 & 8 & \textbf{Etapa 1 e 2 (Planejamento e Identificação):} Início da busca/seleção de estudos e redação do \textbf{Capítulo 3 - Metodologia}. \\ \hline
        Janeiro & 03-01-2026 até 12-01-2026 & 10 & Conclusão da \textbf{Etapa 1 e 2} e da redação do \textbf{Capítulo 3}. \\ \hline
        Janeiro & 13-01-2026 até 31-01-2026 & 19 & \textbf{Etapa 3 e 4 (Extração e Visualização):} Análise dos dados e início da redação do \textbf{Capítulo 4 - Resultados}. \\ \hline
        Fevereiro & 01-02-2026 até 06-02-2026 & 6 & Conclusão da redação do \textbf{Capítulo 4 - Resultados}. \\ \hline
        Fevereiro & 07-02-2026 até 15-02-2026 & 9 & \textbf{Etapa 5 (Relatório):} Redação do \textbf{Capítulo 5 - Considerações Finais} e entrega da versão completa para a orientadora. \\ \hline
        Fevereiro & 16-02-2026 até 28-02-2026 & 13 & Período de revisão e aplicação dos ajustes finais. \\ \hline
        Março & A definir em 03/2026 & - & Preparação da apresentação e \textbf{Defesa do Trabalho de Conclusão de Curso}. \\ \hline
    \end{tabular}
    
    \fonte{Elaboração própria.}
\end{quadro}


% ----------------------------------------------------------
% ELEMENTOS PÓS-TEXTUAIS
% ----------------------------------------------------------
\postextual


%% Referências bibliográficas
\bibliography{referencias.bib}\thispagestyle{empty}
%% ---

% Caso seja necessário glossário em seu documento, fazer a mão ou consulte o manual do pacote abnTeX2 para orientações sobre glossários prontos.

% FUTURE-TODO: implementar modelo simples de glossário
%% Glossário (opcional)
% \begin{glossario}
%     \item[LaTeX] Ferramenta de computador para autoria de documentos criada por D. E. Knuth
%     \item[abnTeX2] Suíte para LaTeX que atende os requisitos das normas da ABNT para elaboração de documentos técnicos e científicos brasileiros
% \end{glossario}
%% ---

% Caso sejam necessários apêndices ou anexos em seu documento use os ambientes abaixo

% %% Apêndices
% \begin{apendicesenv} \thispagestyle{empty}
%     \partapendices % Imprime uma página indicando o início dos apêndices

%     \chapter{Primeiro Apêndice} \thispagestyle{empty}
%     % inserir seu apêndice abaixo

%     \chapter{Segundo Apêndice} \thispagestyle{empty}
%     % inserir seu apêndice abaixo
% \end{apendicesenv}
% %% ---

% %% Anexos
% \begin{anexosenv}
%     \partanexos % Imprime uma página indicando o início dos anexos

%     \chapter{Primeiro Anexo} \thispagestyle{empty}
%     % inserir seu anexo abaixo

%     \chapter{Segundo Anexo} \thispagestyle{empty}
%     % inserir seu anexo abaixo
% \end{anexosenv}
% %% ---


\end{document}